\subsection{Construcción estatal de las doce fibras ciclotómicas}
\label{nuclear:fibras-app-witt}
Sobre las marcas de APP, la multiplicación \(a:r\mapsto4r\) en
\(\ZZ/9\ZZ\) satisface \(a^3=1\). Las hojas verifican
\[
 S(ax,ay)=aS(x,y),\qquad P(ax,ay)=a^{-1}P(x,y),
\]
pues \(16\equiv7\equiv4^{-1}\pmod9\). Las órbitas unitarias
ordenadas son \(\mathcal O_0=(1,4,7)\), \(\mathcal O_1=(2,8,5)\).
En cada una, el módulo de aumento
\(A(\mathcal O)=\ker(\ZZ[\mathcal O]\xrightarrow{\sum}\ZZ)\)
hereda el producto que hace ortonormales los \(e_r\). En la base
\(b_1=e_r-e_{4r}\), \(b_2=e_{4r}-e_{7r}\), sus matrices son
\[
 G_{A_2}=\begin{pmatrix}2&-1\\-1&2\end{pmatrix},\qquad
 C=\begin{pmatrix}0&-1\\1&-1\end{pmatrix},\qquad
 C^{\mathsf T}G_{A_2}C=G_{A_2}.
\]
La forma reticular procede, por tanto, de las órbitas de las marcas.
Para los pares \((1,2),(2,3),(3,1)\) y las hojas
\(\mu\in\{\Sigma,\Pi\}\), se obtiene
\[
 W_{\APP}=\bigoplus_{(i,j)}\bigoplus_\mu\bigoplus_{\varepsilon=0}^1
 A(\mathcal O_\varepsilon),\qquad
 \ell^\Sigma_{ij}(x)=x_i+x_j,\quad
 \ell^\Pi_{ij}(x)=x_ix_j\pmod9.
\]
Cada sumando conserva par, hoja y órbita. La aplicación estatal es
\[
 \Psi:(\ZZ/9\ZZ)^3\longrightarrow W_{\APP},\qquad
 \Psi(x)_{ij,\mu,\varepsilon}=\beta_\varepsilon(\ell^\mu_{ij}(x)),
 \quad
 \beta_\varepsilon(r)=
 \begin{cases}e_r-e_{4r},&r\in\mathcal O_\varepsilon,\\0,&r\notin\mathcal O_\varepsilon.
 \end{cases}
\]
\begin{proposition}[Equivarianza y rango del módulo estatal]
\label{nuclear:psi-rango}
La acción de \(C\) en las seis fibras aditivas y de \(C^{-1}\) en
las seis multiplicativas cumple \(\Psi(ax)=\rho(a)\Psi(x)\).
Las imágenes generan \(W_{\APP}\otimes\QQ\), de dimensión veinticuatro.
\end{proposition}
\begin{proof}
Las identidades de las hojas y \(\beta(4r)=C\beta(r)\) prueban la
equivarianza. Para el rango, se ordenan las coordenadas por los tres
pares indicados, por \(\Sigma,\Pi\), por \(\mathcal O_0,\mathcal O_1\)
y por \(b_1,b_2\). En cada órbita, \(\beta\) toma sucesivamente
\((1,0),(0,1),(-1,-1)\). La matriz de las imágenes de estos estados,
leídos por filas, tiene determinante \(9\):
\[
\begin{array}{rrrr}
(0,0,1)&(0,0,2)&(0,0,4)&(0,0,5)\\
(0,1,0)&(0,1,1)&(0,1,2)&(0,1,3)\\
(0,1,4)&(0,1,5)&(0,1,6)&(0,2,0)\\
(0,2,2)&(0,2,3)&(0,4,0)&(0,5,0)\\
(1,0,1)&(1,0,2)&(1,0,4)&(1,0,5)\\
(1,1,0)&(1,2,0)&(1,4,0)&(1,5,0).
\end{array}
\]
La eliminación sucesiva, reduciendo cada fila por las anteriores y
normalizando su primer coeficiente no anulado, da columnas pivote
numeradas desde cero
\[
\begin{split}
&(8,10,9,11,0,12,14,16,13,15,17,2,\\
&\hspace{2em}18,19,1,3,20,22,21,23,4,6,5,7)
\end{split}
\]
y pivotes
\[
(1,1,1,-1,1,1,1,1,1,-1,3,1,1,3,1,-1,1,1,1,-1,1,1,1,-1).
\]
Su producto, multiplicado por el signo de la permutación de columnas,
es \(9\). El menor invertible acredita el rango total.
\end{proof}

\subsection{Alineamiento dodecafásico con las posiciones de Witt}
\label{nuclear:alineamiento-witt}
El índice del par es \(i\in\ZZ/3\ZZ\), el de la hoja
\(\mu\in\{0,1\}\) y el de la órbita \(\varepsilon\in\{0,1\}\).
Fijado el origen dodecafásico, \(\lambda(i,\mu,\varepsilon)
=4i+2\mu+\varepsilon\pmod{12}\) es biyectiva: los cuatro valores
de cada \(i\) forman un bloque y los tres bloques son disjuntos.
Sea \(R=\left(\begin{smallmatrix}0&1\\1&0\end{smallmatrix}\right)\).
La multiplicación prueba \(RCR=C^{-1}\) y
\(R^{\mathsf T}G_{A_2}R=G_{A_2}\). Para las trivializaciones
\(\iota_b\) y los marcos \(P_b\) transportados conjuntamente con
la incidencia de Witt, ponemos
\[
g_b=\iota_b^{-1}P_bCP_b^{-1}\iota_b,\qquad
T_{i\mu\varepsilon}=\iota_{\lambda(i,\mu,\varepsilon)}^{-1}
P_{\lambda(i,\mu,\varepsilon)}R^\mu.
\]
Entonces
\[
g_{\lambda(i,\mu,\varepsilon)}T_{i\mu\varepsilon}
=T_{i\mu\varepsilon}C^{(-1)^\mu}.
\]
La suma directa es un isomorfismo \(C_3\)-equivariante entre las doce
fibras estatales y las doce posiciones reticulares. Conserva par,
hoja, órbita y forma bajo el transporte simultáneo del origen,
orden y orientación. Estos datos permanecen en el pegado reticular
y el paso al vecino que se construyen a continuación.
